\begin{filecontents*}{\jobname.bib}
@inproceedings{pourreza2023dinsql,
  author    = {Mohammadreza Pourreza and Davood Rafiei},
  title     = {{DIN-SQL}: Decomposed In-Context Learning of Text-to-{SQL} with Self-Correction},
  booktitle = {Advances in Neural Information Processing Systems (NeurIPS)},
  volume    = {36},
  pages     = {37269--37286},
  year      = {2023}
}

@article{gao2024dailsql,
  author    = {Dawei Gao and Haibin Wang and Yaliang Li and Xiuyu Sun and Yichen Qian and Bolin Ding and Jingren Zhou},
  title     = {Text-to-{SQL} Empowered by Large Language Models: A Benchmark Evaluation},
  journal   = {Proceedings of the VLDB Endowment},
  volume    = {17},
  number    = {5},
  pages     = {1132--1145},
  year      = {2024},
  doi       = {10.14778/3641204.3641221}
}

@article{qwen2024qwen25,
  author    = {{Qwen Team}},
  title     = {{Qwen2.5} Technical Report},
  journal   = {arXiv preprint arXiv:2412.15115},
  year      = {2024}
}

@inproceedings{madaan2023selfrefine,
  author    = {Aman Madaan and Niket Tandon and Prakhar Gupta and Skyler Hallinan and Luyu Gao and Sarah Wiegreffe and Uri Alon and Nouha Dziri and Shrimai Prabhumoye and Yiming Yang and Shashank Gupta and Bodhisattwa Prasad Majumder and Katherine Hermann and Sean Welleck and Amir Yazdanbakhsh and Peter Clark},
  title     = {Self-Refine: Iterative Refinement with Self-Feedback},
  booktitle = {Advances in Neural Information Processing Systems (NeurIPS)},
  volume    = {36},
  pages     = {46534--46594},
  year      = {2023}
}

@inproceedings{shinn2023reflexion,
  author    = {Noah Shinn and Federico Cassano and Ashwin Gopinath and Karthik Narasimhan and Shunyu Yao},
  title     = {Reflexion: Language Agents with Verbal Reinforcement Learning},
  booktitle = {Advances in Neural Information Processing Systems (NeurIPS)},
  volume    = {36},
  pages     = {8634--8652},
  year      = {2023}
}

@inproceedings{lewis2020rag,
  author    = {Patrick Lewis and Ethan Perez and Aleksandra Piktus and Fabio Petroni and Vladimir Karpukhin and Naman Goyal and Heinrich K{\"u}ttler and Mike Lewis and Wen-tau Yih and Tim Rockt{\"a}schel and Sebastian Riedel and Douwe Kiela},
  title     = {Retrieval-Augmented Generation for Knowledge-Intensive {NLP} Tasks},
  booktitle = {Advances in Neural Information Processing Systems (NeurIPS)},
  volume    = {33},
  pages     = {9459--9474},
  year      = {2020}
}

@inproceedings{park2023generativeagents,
  author    = {Joon Sung Park and Joseph C. O'Brien and Carrie J. Cai and Meredith Ringel Morris and Percy Liang and Michael S. Bernstein},
  title     = {Generative Agents: Interactive Simulacra of Human Behavior},
  booktitle = {Proceedings of the 36th Annual ACM Symposium on User Interface Software and Technology (UIST)},
  pages     = {1--22},
  year      = {2023},
  doi       = {10.1145/3586183.3606763}
}

@article{packer2023memgpt,
  author    = {Charles Packer and Vivian Fang and Shishir G. Patil and Kevin Lin and Sarah Wooders and Joseph E. Gonzalez},
  title     = {{MemGPT}: Towards {LLMs} as Operating Systems},
  journal   = {arXiv preprint arXiv:2310.08560},
  year      = {2023}
}

@inproceedings{rebedea2023nemoguardrails,
  author    = {Traian Rebedea and Razvan Dinu and Makesh Narsimhan Sreedhar and Christopher Parisien and Jonathan Cohen},
  title     = {{NeMo} Guardrails: A Toolkit for Controllable and Safe {LLM} Applications with Programmable Rails},
  booktitle = {Proceedings of the 2023 Conference on Empirical Methods in Natural Language Processing (EMNLP): System Demonstrations},
  pages     = {431--445},
  year      = {2023}
}

@misc{mao2023sqlglot,
  author       = {Toby Mao},
  title        = {{SQLGlot}: An Extensible {SQL} Parser and Transpiler},
  year         = {2023},
  howpublished = {\url{https://github.com/tobymao/sqlglot}}
}

@inproceedings{halfond2005amnesia,
  author    = {William G. J. Halfond and Alessandro Orso},
  title     = {{AMNESIA}: Analysis and Monitoring for Neutralizing {SQL}-Injection Attacks},
  booktitle = {Proceedings of the 20th IEEE/ACM International Conference on Automated Software Engineering (ASE)},
  pages     = {174--183},
  year      = {2005},
  doi       = {10.1145/1101908.1101935}
}

@inproceedings{finegandollak2018eval,
  author    = {Catherine Finegan-Dollak and Jonathan K. Kummerfeld and Li Zhang and Karthik Ramanathan and Sesh Sadasivam and Rui Zhang and Dragomir Radev},
  title     = {Improving Text-to-{SQL} Evaluation Methodology},
  booktitle = {Proceedings of the 56th Annual Meeting of the Association for Computational Linguistics (ACL)},
  pages     = {351--360},
  year      = {2018},
  doi       = {10.18653/v1/P18-1033}
}

@inproceedings{zhong2020distilled,
  author    = {Ruiqi Zhong and Tao Yu and Dan Klein},
  title     = {Semantic Evaluation for Text-to-{SQL} with Distilled Test Suites},
  booktitle = {Proceedings of the 2020 Conference on Empirical Methods in Natural Language Processing (EMNLP)},
  pages     = {396--411},
  year      = {2020},
  doi       = {10.18653/v1/2020.emnlp-main.29}
}

@inproceedings{wang2025macsql,
  author    = {Bing Wang and Changyu Ren and Jian Yang and Xinnian Liang and Jiaqi Bai and Linzheng Chai and Zhao Yan and Qian-Wen Zhang and Di Yin and Xing Sun and Zhoujun Li},
  title     = {{MAC-SQL}: A Multi-Agent Collaborative Framework for Text-to-{SQL}},
  booktitle = {Proceedings of the 31st International Conference on Computational Linguistics (COLING)},
  pages     = {1--15},
  year      = {2025}
}

@inproceedings{chen2024selfdebug,
  author    = {Xinyun Chen and Maxwell Lin and Nathanael Sch{\"a}rli and Denny Zhou},
  title     = {Teaching Large Language Models to Self-Debug},
  booktitle = {Proceedings of the International Conference on Learning Representations (ICLR)},
  year      = {2024}
}

@inproceedings{zhang2023selfedit,
  author    = {Kechi Zhang and Zhuo Li and Jia Li and Ge Li and Zhi Jin},
  title     = {Self-Edit: Fault-Aware Code Editor for Code Generation},
  booktitle = {Proceedings of the 61st Annual Meeting of the Association for Computational Linguistics (ACL)},
  pages     = {769--787},
  year      = {2023},
  doi       = {10.18653/v1/2023.acl-long.45}
}

@article{johnson2021faiss,
  author    = {Jeff Johnson and Matthijs Douze and Herv{\'e} J{\'e}gou},
  title     = {Billion-Scale Similarity Search with {GPUs}},
  journal   = {IEEE Transactions on Big Data},
  volume    = {7},
  number    = {3},
  pages     = {535--547},
  year      = {2021},
  doi       = {10.1109/TBDATA.2019.2921572}
}

@article{nussbaum2024nomic,
  author    = {Zach Nussbaum and John X. Morris and Brandon Duderstadt and Andriy Mulyar},
  title     = {Nomic Embed: Training a Reproducible Long Context Text Embedder},
  journal   = {arXiv preprint arXiv:2402.01613},
  year      = {2024}
}

@article{sumers2024coala,
  author    = {Theodore Sumers and Shunyu Yao and Karthik Narasimhan and Thomas L. Griffiths},
  title     = {Cognitive Architectures for Language Agents},
  journal   = {Transactions on Machine Learning Research (TMLR)},
  year      = {2024}
}

@inproceedings{wei2023jailbroken,
  author    = {Alexander Wei and Nika Haghtalab and Jacob Steinhardt},
  title     = {Jailbroken: How Does {LLM} Safety Training Fail?},
  booktitle = {Advances in Neural Information Processing Systems (NeurIPS)},
  volume    = {36},
  pages     = {80079--80110},
  year      = {2023}
}

@article{zou2023universal,
  author    = {Andy Zou and Zifan Wang and Nicholas Carlini and Milad Nasr and J. Zico Kolter and Matt Fredrikson},
  title     = {Universal and Transferable Adversarial Attacks on Aligned Language Models},
  journal   = {arXiv preprint arXiv:2307.15043},
  year      = {2023}
}

@misc{langgraph2024,
  author       = {{LangChain}},
  title        = {{LangGraph}: Build Resilient Language Agents as Graphs},
  year         = {2024},
  howpublished = {\url{https://github.com/langchain-ai/langgraph}}
}

@inproceedings{yu2018spider,
  author    = {Tao Yu and Rui Zhang and Kai Yang and Michihiro Yasunaga and Dongxu Wang and Zifan Li and James Ma and Irene Li and Qingning Yao and Shanelle Roman and Zilin Zhang and Dragomir Radev},
  title     = {Spider: A Large-Scale Human-Labeled Dataset for Complex and Cross-Domain Semantic Parsing and Text-to-{SQL} Task},
  booktitle = {Proceedings of the 2018 Conference on Empirical Methods in Natural Language Processing (EMNLP)},
  pages     = {387--399},
  year      = {2018},
  doi       = {10.18653/v1/D18-1425}
}

@inproceedings{li2023bird,
  author    = {Jinyang Li and Binyuan Hui and Ge Qu and Jiaxi Yang and Binhua Li and Bowen Li and Bailin Wang and Bowen Qin and Ruiying Geng and Nan Huo and Xuanhe Zhou and Chenhao Ma and Guoliang Li and Kevin C. C. Chang and Fei Huang and Reynold Cheng and Yongbin Li},
  title     = {Can {LLM} Already Serve as a Database Interface? {A} {BIg} Bench for Large-Scale Database Grounded Text-to-{SQL}s},
  booktitle = {Advances in Neural Information Processing Systems (NeurIPS)},
  volume    = {36},
  pages     = {64082--64101},
  year      = {2023}
}

@article{wang2018executionguided,
  author    = {Chenglong Wang and Kedar Tatwawadi and Marc Brockschmidt and Po-Sen Huang and Yi Mao and Oleksandr Polozov and Rishabh Singh},
  title     = {Robust Text-to-{SQL} Generation with Execution-Guided Decoding},
  journal   = {arXiv preprint arXiv:1807.03100},
  year      = {2018}
}

@article{wei2022emergent,
  author    = {Jason Wei and Yi Tay and Rishi Bommasani and Colin Raffel and Barret Zoph and Sebastian Borgeaud and Dani Yogatama and Maarten Bosma and Denny Zhou and Donald Metzler and Ed H. Chi and Tatsunori Hashimoto and Oriol Vinyals and Percy Liang and Jeff Dean and William Fedus},
  title     = {Emergent Abilities of Large Language Models},
  journal   = {Transactions on Machine Learning Research (TMLR)},
  year      = {2022}
}
\end{filecontents*}

\documentclass[conference]{IEEEtran}
\IEEEoverridecommandlockouts

\usepackage{amsmath,amssymb,amsfonts}
\usepackage{graphicx}
\usepackage{booktabs}
\usepackage{tabularx}
\usepackage{cite}
\usepackage{url}
\usepackage{centernot}
\usepackage{microtype}

\graphicspath{{figures/}{figures/png/}{figures/svg/}{./}}

% Calibrated line-breaking tolerance
\emergencystretch=3em

% Safe url and nolinkurl handling
\providecommand{\nolinkurl}[1]{\url{#1}}

% Calibrated IEEE Float & Gap Parameters
\setlength{\floatsep}{8.0pt plus 1.0pt minus 1.0pt}
\setlength{\textfloatsep}{10.0pt plus 1.0pt minus 2.0pt}
\setlength{\intextsep}{8.0pt plus 1.0pt minus 1.0pt}
\setlength{\dblfloatsep}{8.0pt plus 1.0pt minus 1.0pt}
\setlength{\dbltextfloatsep}{10.0pt plus 1.0pt minus 2.0pt}

% Calibrated Display Math Spacing
\setlength{\abovedisplayskip}{4.5pt plus 1.0pt minus 1.0pt}
\setlength{\belowdisplayskip}{4.5pt plus 1.0pt minus 1.0pt}
\setlength{\abovedisplayshortskip}{2.0pt plus 1.0pt}
\setlength{\belowdisplayshortskip}{2.0pt plus 1.0pt}

% Section, Subsection, and Subsubsection Breathing Space Calibration
\makeatletter
\def\section{\@startsection{section}{1}{\z@}{2.8ex plus 1.2ex minus 0.5ex}{1.5ex plus 0.4ex minus 0.2ex}{\normalfont\normalsize\centering\scshape}}
\def\subsection{\@startsection{subsection}{2}{\z@}{2.2ex plus 1.0ex minus 0.4ex}{1.0ex plus 0.3ex minus 0.2ex}{\normalfont\normalsize\itshape}}
\def\subsubsection{\@startsection{subsubsection}{3}{\z@}{1.8ex plus 0.8ex minus 0.3ex}{0.8ex plus 0.2ex minus 0.1ex}{\normalfont\normalsize\itshape}}

% Flawless Table and Figure Caption Formatting: Eliminates Overlapping Lines and Preserves Clean Spacing
\long\def\@makecaption#1#2{%
  \ifx\@captype\@IEEEtablestring
    \setbox\@tempboxa\hbox{\footnotesize #1}%
    \ifdim \wd\@tempboxa >\hsize
      {\footnotesize #1\par}%
    \else
      \hbox to\hsize{\footnotesize\hfil\box\@tempboxa\hfil}%
    \fi
    \vskip 2.5pt
    \setbox\@tempboxa\hbox{\footnotesize\scshape #2}%
    \ifdim \wd\@tempboxa >\hsize
      {\footnotesize\scshape #2\par}%
    \else
      \hbox to\hsize{\footnotesize\hfil\box\@tempboxa\hfil}%
    \fi
    \vskip 6.0pt plus 1.0pt minus 1.0pt % Clean breathing space before \toprule
  \else
    \@IEEEfigurecaptionsepspace
    \setbox\@tempboxa\hbox{\footnotesize\itshape #2}%
    \ifdim \wd\@tempboxa >\hsize
      {\footnotesize\itshape #2\par}%
    \else
      \hbox to\hsize{\footnotesize\itshape\hfil\box\@tempboxa\hfil}%
    \fi
    \vskip 4.0pt plus 1.0pt minus 1.0pt
  \fi}
\makeatother

\begin{document}

\title{Adaptive Runtime Memory Governance: Closed-Loop Operational Knowledge Extraction, Repair Efficiency, and Execution Safety in Local Text-to-SQL Systems}

% Symmetrical V-Shape (Inverted Pyramid) Author Architecture
\author{
  % Row 1: 2 Authors (Wide apart)
  \makebox[\textwidth][c]{%
    \begin{minipage}[t]{0.36\textwidth}
      \centering
      \textbf{Inampudi Govardhana Rao}\\
      \textit{School of Computer Science and Engineering}\\
      \footnotesize govardhanarao.i@vitap.ac.in
    \end{minipage}%
    \hspace{0.14\textwidth}%
    \begin{minipage}[t]{0.36\textwidth}
      \centering
      \textbf{Ghanta Sundar Siddhartha}\\
      \textit{B.Tech CSE Core}\\
      \footnotesize siddharthaghanta10@gmail.com
    \end{minipage}%
  }\\[1.5ex]
  % Row 2: 2 Authors (Inward)
  \makebox[\textwidth][c]{%
    \begin{minipage}[t]{0.38\textwidth}
      \centering
      \textbf{Sudarsanan Shilpa}\\
      \textit{B.Tech CSE Core}\\
      \footnotesize shilpasudarsanan4@gmail.com
    \end{minipage}%
    \hspace{0.06\textwidth}%
    \begin{minipage}[t]{0.38\textwidth}
      \centering
      \textbf{Bommadevara S N V Datta Prasada Rayulu}\\
      \textit{B.Tech CSE Core}\\
      \footnotesize dattabommadevara123@gmail.com
    \end{minipage}%
  }\\[1.5ex]
  % Row 3: 1 Author (V Apex Center)
  \makebox[\textwidth][c]{%
    \begin{minipage}[t]{0.36\textwidth}
      \centering
      \textbf{Velpuri Danaiah}\\
      \textit{B.Tech CSE Core}\\
      \footnotesize danaiahvelpuri78@gmail.com
    \end{minipage}%
  }
}

\maketitle

\begin{abstract}
Deploying local Large Language Models (LLMs) for enterprise Text-to-SQL workflows introduces acute operational challenges: repetitive schema repair oscillations, unconstrained vector memory accumulation in retrieval stores, and the catastrophic risk of executing destructive SQL mutations against relational warehouses. Standard stateless self-correction prompts models with raw database errors without cross-query operational memory, while unmanaged vector retrieval systems suffer from duplicate accumulation and memory bloat. This paper presents \textbf{Adaptive Runtime Memory Governance (ARMG)}, a closed-loop runtime architecture that converts PostgreSQL database execution feedback into governed, reusable operational memory. ARMG wraps frozen local open-weights models in a 10-node directed state graph integrating deterministic Abstract Syntax Tree (AST) and catalog-based error diagnosis across a canonical 7-tier exception taxonomy, ephemeral knowledge extraction, admission- and reinforcement-gated FAISS vector memory, diagnostic negative constraints, and a deterministic pre-execution AST safety guardrail.

In an empirical evaluation across three isolated repeated benchmark executions ($n = 3$, comprising 450 total evaluated query runs across six experimental modes) on a 25-query synthetic enterprise-style Star Schema data warehouse using local \texttt{qwen2.5:7b-instruct} under greedy decoding (\texttt{temperature = 0.0}), ARMG enforced pre-execution safety: across 450 post-remediation evaluations and 150 historical baseline evaluations (600 audited executions in total), no destructive SQL statement reached PostgreSQL. Relative to stateless self-correction, ARMG achieved an observed \textbf{34.38\% reduction in mean repair iterations} (0.28 vs. 0.43 retries per query) and a \textbf{5.30\% reduction in token consumption} (542.37 vs. 572.72 tokens per query), while increasing PostgreSQL execution success from 92.00\% to 96.00\%. Furthermore, ARMG enforces algorithmic mutual exclusion between reinforcement and new admission; in the evaluated sequential benchmark, the Mode 4 store remained at exactly 3 memories from Q15 to Q25, whereas the unmanaged Naive Vector RAG control accumulated 23 entries. However, in the evaluated benchmark, ARMG exhibited a \textbf{19.79\% higher end-to-end latency} (8,564.89~ms vs. 7,149.68~ms) while using fewer mean repair iterations, and \textbf{relational execution accuracy remained at 68.00\%} (17 of 25 queries) across both ARMG and stateless self-correction. In addition, continuous exponential temporal decay remained experimentally unexercised under the $\sim$3.5-minute benchmark execution clock. These findings demonstrate that governed operational memory provides bounded repair behavior, token economy, and pre-execution safety enforcement in the evaluated setting, while establishing that operational memory did not overcome the observed semantic errors on the evaluated window-function queries in the local 7B configuration.
\end{abstract}

\begin{IEEEkeywords}
Text-to-SQL, Large Language Models, Runtime Memory Governance, Vector Retrieval, Database Safety, Query Repair, LangGraph, Enterprise Data Warehouses.
\end{IEEEkeywords}

\section{Introduction}
Enterprise adoption of natural language interfaces to relational databases (Text-to-SQL) has accelerated with advances in instruction-tuned Large Language Models \cite{pourreza2023dinsql, gao2024dailsql}. However, organizations subject to stringent data privacy, sovereignty, or computational cost constraints increasingly mandate on-premises deployment using localized open-weights models (e.g., 7B-parameter architectures) \cite{qwen2024qwen25}. Deploying smaller LLMs locally introduces critical operational challenges:
\begin{enumerate}
    \item \textbf{Repetitive Repair Failures}: Stateless self-correction approaches \cite{madaan2023selfrefine, shinn2023reflexion} prompt the language model with raw database driver execution exceptions. Without persistent cross-query memory, the model frequently oscillates between identical invalid SQL syntax or schema constructs across successive repair cycles.
    \item \textbf{Retrieval Corruption \& Memory Bloat}: Naive implementations of retrieval augmentation \cite{lewis2020rag} that append historical execution exemplars without explicit lifecycle governance risk accumulating duplicate, conflicting, or stale entries over time \cite{park2023generativeagents, packer2023memgpt}.
    \item \textbf{Execution Safety Hazards}: Unbounded Text-to-SQL agents risk synthesizing destructive Data Manipulation Language (DML) or Data Definition Language (DDL) statements (\texttt{DELETE}, \allowbreak\texttt{DROP}, \allowbreak\texttt{UPDATE}, \allowbreak\texttt{TRUNCATE}), presenting severe operational risks if connected to live warehouse environments \cite{rebedea2023nemoguardrails, mao2023sqlglot, halfond2005amnesia}.
    \item \textbf{The Executable vs. Relational Correctness Gap}: While commercial evaluations frequently report database execution rates, executable SQL is frequently non-equivalent to the analytical intent of the user \cite{finegandollak2018eval, zhong2020distilled}.
\end{enumerate}

To address these limitations, we introduce \textbf{Adaptive Runtime Memory Governance (ARMG)}, a closed-loop runtime architecture that governs the extraction, admission, retrieval, reinforcement, and application of runtime operational knowledge for local Text-to-SQL systems. Rather than updating underlying model weights via fine-tuning, ARMG structures execution feedback into an ephemeral diagnostic intermediate representation (\texttt{Runtime\-Knowledge}) that undergoes mathematical admission gating before entering a persistent FAISS vector index. Crucially, ARMG enforces strict algorithmic mutual exclusion between the reinforcement of existing memories and the admission of new knowledge, preventing duplicate memory accumulation. A deterministic Abstract Syntax Tree (AST) validation layer parses every query prior to database driver invocation, halting destructive mutations before execution.

\subsection{Research Gap Formulation}
The reviewed literature emphasizes large-model Text-to-SQL generation, self-correction, multi-agent workflows, and benchmark-level evaluation; fewer works address persistent runtime memory governance in local deployments. Within the literature reviewed in this study, we did not identify a framework that jointly evaluates persistent operational memory governance, runtime repair, and deterministic pre-execution AST safety in the evaluated local Text-to-SQL setting. \textbf{ARMG investigates an integrated runtime governance architecture combining:}
\begin{enumerate}
    \item Local open-weights Text-to-SQL inference;
    \item Runtime physical database execution feedback;
    \item Zero-token deterministic 7-tier exception diagnosis;
    \item Ephemeral intermediate knowledge representation;
    \item Governed persistent operational vector memory;
    \item Algorithmic mutual exclusion between memory reinforcement and admission;
    \item Diagnostic negative repair constraints;
    \item Deterministic pre-execution SQL AST safety validation; and
    \item Rigorous, explicit evaluation separating PostgreSQL execution success from relational execution accuracy.
\end{enumerate}

\subsection{Summary of Empirical Findings}
We evaluate ARMG across three isolated repeated benchmark executions ($n = 3$, comprising 450 total evaluated query runs across six experimental modes) on a synthetically generated enterprise-style Star Schema data warehouse benchmark. Our experimental findings establish:
\begin{itemize}
    \item \textbf{Verified Retrieval Geometry}: Unit-$L_2$ normalization of runtime query embeddings restored the intended FAISS $L_2$ nearest-neighbor similarity geometry above retrieval threshold $\tau = 0.50$, yielding 16 retrieval events across 12 distinct benchmark queries ($48.0\%$ query coverage).
    \item \textbf{Controlled Lifecycle \& Deduplication}: In the evaluated sequential benchmark, algorithmic mutual exclusion between reinforcement and new admission maintained persistent vector store size at 3 memories across queries Q15 through Q25, whereas the evaluated Naive Vector RAG control accumulated 23 entries.
    \item \textbf{Repair-Loop Efficiency \& Token Economy}: ARMG achieved an observed 34.38\% reduction in mean repair loop iterations (0.28 vs. 0.43 retries per query) and a 5.30\% reduction in mean token expenditure (542.37 vs. 572.72 tokens) compared to stateless self-correction.
    \item \textbf{Verified Pre-Execution Safety}: Across 450 post-remediation evaluations and 150 historical baseline evaluations (600 audited executions in total), no destructive SQL statement reached PostgreSQL in the evaluated benchmark.
    \item \textbf{Architectural Latency Trade-Off}: In the evaluated benchmark, ARMG exhibited a 19.79\% higher end-to-end latency (8,564.89~ms vs. 7,149.68~ms) while using fewer mean repair iterations, which coincided with the additional embedding, retrieval, and state-graph processing stages.
    \item \textbf{Relational Execution Accuracy Plateau}: Relational execution accuracy remained identical at 68.00\% (17/25 queries) between Full ARMG and stateless self-correction; in the evaluated benchmark, operational memory did not overcome the observed semantic errors on the more complex window-function queries.
\end{itemize}

\subsection{Research Contributions}
\begin{enumerate}
    \item \textbf{Architectural Contribution}: We design and implement a closed-loop runtime architecture decoupled from model weights, orchestrating zero-token catalog schema introspection, static AST safety validation, deterministic error diagnosis without LLM inference, and bounded self-correction in a 10-node directed state graph.
    \item \textbf{Runtime Operational Knowledge Contribution}: We formalize an immutable ephemeral knowledge artifact (\texttt{Runtime\-Knowledge}) and a diagnostic negative constraint mechanism that isolates broken SQL tokens and schema identifiers, preventing repetitive repair cycling.
    \item \textbf{Memory Governance Contribution}: We formulate a mathematical governance engine incorporating multi-factor operational utility, asymptotic confidence escalation, multiplicative failure penalties, continuous exponential decay, and an algorithmic mutual-exclusion deduplication invariant.
    \item \textbf{Safety Contribution}: We implement a multi-layered pre-execution AST containment mechanism using SQLGlot that intercepts destructive DDL/DML mutations and multi-statement injections prior to database driver invocation.
    \item \textbf{Empirical Characterization of Local 7B Text-to-SQL}: We provide a rigorous evaluation across 450 query runs, isolating the trade-off profile between repair efficiency, token expenditure, latency overhead, and the observed semantic reasoning boundaries of the local 7B configuration.
\end{enumerate}

\section{Related Work}

\subsection{Text-to-SQL and Self-Correction}
Modern Text-to-SQL research has progressed from specialized sequence-to-sequence architectures and relation-aware graph encoders to multi-stage in-context learning pipelines utilizing decomposed reasoning, schema pruning, and execution-guided self-correction \cite{pourreza2023dinsql, gao2024dailsql, wang2025macsql, wang2018executionguided}. Decomposed In-Context Learning (DIN-SQL) \cite{pourreza2023dinsql} breaks the generation task into sub-tasks (schema linking, classification, SQL generation, and self-correction), achieving substantial improvements on complex benchmarks. Similarly, DAIL-SQL \cite{gao2024dailsql} systematically benchmarks prompt representation and selection strategies, demonstrating that token efficiency is critical for effective few-shot prompting. Multi-agent collaborative frameworks, such as MAC-SQL \cite{wang2025macsql}, deploy specialized agents (Decomposer, Selector, Refiner) to coordinate complex analytical reasoning.

In parallel, execution-guided decoding and repair frameworks leverage runtime feedback to correct invalid SQL \cite{madaan2023selfrefine, shinn2023reflexion, chen2024selfdebug, zhang2023selfedit, wang2018executionguided}. Stateless self-correction approaches like Self-Refine \cite{madaan2023selfrefine} and Reflexion \cite{shinn2023reflexion} utilize iterative prompt feedback loops. Self-Debug \cite{chen2024selfdebug} and Self-Edit \cite{zhang2023selfedit} integrate error traces and test-case execution results to guide model self-repair. However, these systems operate exclusively within the context of an individual query: once query synthesis completes, the repair trace is discarded. Consequently, stateless models cannot transfer operational repair knowledge across sequential queries in a multi-query session, leading to repeated repair failures on recurring schema traps. In contrast, ARMG extracts persistent, cross-query operational knowledge and injects deterministic negative constraints to break repair oscillation loops.

\subsection{Retrieval-Augmented Generation and Agent Memory}
Retrieval-Augmented Generation (RAG) frameworks augment parametric language model weights with non-parametric dense vector indices \cite{lewis2020rag}, typically indexed via high-performance nearest-neighbor search libraries such as FAISS \cite{johnson2021faiss} and dense embedding models \cite{nussbaum2024nomic}. In agentic workflows, memory streams have been introduced to maintain long-term behavioral consistency and historical context \cite{park2023generativeagents, packer2023memgpt, sumers2024coala}. Generative Agents \cite{park2023generativeagents} introduced mathematical scoring heuristics based on recency, importance, and relevance, including exponential decay. MemGPT \cite{packer2023memgpt} formalized hierarchical virtual memory management to mitigate bounded LLM context windows, while the CoALA cognitive architecture \cite{sumers2024coala} systematically categorized working, episodic, and semantic memory in language agents.

While these memory-augmented frameworks establish foundational principles for agent behavior, they focus primarily on conversational dialogue or general decision tasks rather than the operational constraints of database execution. Furthermore, in naive retrieval implementations that append historical execution exemplars without explicit lifecycle governance, stores can accumulate redundant or conflicting entries. In contrast, ARMG establishes an explicit operational lifecycle:
\begin{equation}
\begin{aligned}
\text{Observation} &\to \text{Diagnosis} \to \text{Knowledge} \\
&\to \text{Governance} \to \text{Memory}
\end{aligned}
\end{equation}
ARMG explicitly distinguishes between \texttt{Runtime\-Knowledge}---an immutable, ephemeral diagnostic artifact produced during query repair---and \texttt{Runtime\-Memory}---a persistent, mathematically governed operational memory record admitted to FAISS vector storage only upon passing strict utility thresholds and mutual-exclusion deduplication checks. ARMG evaluates candidate knowledge against an admission threshold ($\text{Utility}_0 \ge 0.25$) and enforces strict algorithmic mutual exclusion between memory reinforcement and new admission, maintaining persistent vector store size strictly bounded.

\subsection{Database Safety and Guardrails}
Prompt-based safety directives instruct LLMs to avoid destructive commands via natural language system prompts. However, extensive empirical literature demonstrates that prompt-based alignment is fundamentally vulnerable to adversarial jailbreaks, semantic confusion, and prompt-injection attacks \cite{wei2023jailbroken, zou2023universal}. Programmable safety frameworks, such as NeMo Guardrails \cite{rebedea2023nemoguardrails}, attempt to guide conversational paths, but remain heuristic and probabilistic when applied to executable database code.

In database security, deterministic validation via Abstract Syntax Tree (AST) analysis provides deterministic execution boundaries that decouple policy enforcement from model behavior \cite{mao2023sqlglot, halfond2005amnesia}. Classic systems like AMNESIA \cite{halfond2005amnesia} established that comparing runtime SQL queries against static syntactic models prevents injection attacks. ARMG builds upon this principle by embedding static AST parsing via SQLGlot \cite{mao2023sqlglot} directly into the runtime state graph. ARMG's safety guard acts as a deterministic pre-execution containment mechanism, inspecting statement counts, AST root nodes, and expression types to ensure that destructive DDL/DML mutations are halted before database driver invocation.

\section{Problem Formulation}

Let an enterprise relational data warehouse schema be defined as a tuple:
\begin{equation}
\mathcal{S} = (\mathcal{T}, \mathcal{C}, \mathcal{R})
\end{equation}
where $\mathcal{T} = \{T_1, T_2, \dots, T_m\}$ represents the set of relational tables, $\mathcal{C} = \{c_{i,1}, c_{i,2}, \dots, c_{i,k}\}$ represents the set of typed columns for table $T_i$, and $\mathcal{R} = \{(c_{i,a}, c_{j,b})\}$ represents foreign-key integrity constraints.

\subsection{SQL Generation Task}
Given a natural language analytical question $q \in \mathcal{Q}$ and a deterministically pruned schema context $\mathcal{S}_q \subseteq \mathcal{S}$, an initial generator $\mathcal{G}_{\theta}$ parameterized by frozen local model weights $\theta$ synthesizes candidate query $s_0$:
\begin{equation}
s_0 = \mathcal{G}_{\theta}(q, \mathcal{S}_q, \mathcal{M}_{\text{ret}})
\end{equation}
where $\mathcal{M}_{\text{ret}}$ denotes operational memories retrieved from persistent storage.

\subsection{Pre-Execution Static Safety Validation}
Before database submission, candidate query $s_k$ passes through deterministic AST validation $\mathcal{V}_{\text{AST}}(s_k)$:
\begin{equation}
\begin{aligned}
\mathcal{V}_{\text{AST}}(s_k) \to ( & \text{is\_valid} \in \{\text{True}, \text{False}\}, \\
& \text{error\_reason} )
\end{aligned}
\end{equation}
If $s_k$ contains destructive DDL/DML mutations or non-SELECT root expressions, the execution halts immediately in state $\text{STATUS\_BLOCKED}$, with zero physical database interaction.

\subsection{Execution Environment Feedback}
A validated query is submitted to physical database execution environment $\mathcal{E}$:
\begin{equation}
R_k = \mathcal{E}(s_k) = (\text{status}_k, \mathcal{D}_{s_k}, e_k, t_k)
\end{equation}
where $\text{status}_k \in \{\text{SUCCESS}, \text{FAILURE}\}$, $\mathcal{D}_{s_k}$ is the resulting tuple set, $e_k$ is the raw driver error trace string, and $t_k$ is execution latency.

\subsection{Bounded Runtime Repair Task}
If $\text{status}_k = \text{FAILURE}$, a deterministic diagnostic engine parses normalized error trace $\bar{e}_k$ against schema $\mathcal{S}$ to produce diagnostic tuple $\Delta_k = (\text{taxonomy}, \text{root\_cause}, \mathcal{C}_{\text{cand}}, \mathcal{N}_k, \text{rule})$. A repair agent $\mathcal{R}_{\theta}$ synthesizes replacement query $s_{k+1}$:
\begin{equation}
s_{k+1} = \mathcal{R}_{\theta}(q, \mathcal{S}_q, s_k, \bar{e}_k, \Delta_k, \mathcal{M}_{\text{ret}})
\end{equation}
subject to strict repair constraints forbidding identifiers in negative constraints $\mathcal{N}_k$, bounded by maximum retry budget $K_{\max} = 3$ (total attempts $\le 4$).

\subsection{Distinction: PostgreSQL Execution Success vs. Relational Accuracy}
We formalize two distinct evaluation metrics:
\begin{itemize}
    \item \textbf{PostgreSQL Execution Success ($\text{ExecSucc}$)}: A binary indicator evaluating whether query $s$ executed against PostgreSQL without syntax, schema, or driver exceptions, returning a valid tuple set ($\text{status} = \text{SUCCESS}$). Executability is a necessary but insufficient condition for correctness.
    \item \textbf{Relational Execution Accuracy ($\text{ExecAcc}$)}: A binary indicator evaluated via deterministic relational equivalence comparator $\mathcal{D}_s \equiv_{\text{rel}} \mathcal{D}_{s^*}$ against the ground-truth gold SQL result set $\mathcal{D}_{s^*}$. Evaluates multiset tuple matching, row multiplicity, column attribute alignment, and numerical precision tolerance:
\end{itemize}
\begin{equation}
\text{ExecAcc}(s) = 1 \implies \text{ExecSucc}(s) = 1
\end{equation}
\begin{equation}
\text{ExecSucc}(s) = 1 \centernot\implies \text{ExecAcc}(s) = 1
\end{equation}

\section{ARMG Architecture}
The ARMG architecture is implemented as a closed-loop directed execution graph comprising 10 functional nodes orchestrated via LangGraph \cite{langgraph2024} (\url{graph/workflow.py}):
\begin{enumerate}
    \item \textbf{Node 1: \texttt{introspect\_\allowbreak and\_\allowbreak prune\_\allowbreak node}}: Introspects relational database catalog metadata via zero-token \texttt{information\_schema} queries. Deterministically prunes schema down to query-relevant tables using rule-based token matching.
    \item \textbf{Node 2: \texttt{memory\_\allowbreak retrieval\_\allowbreak node}}: Computes unit-$L_2$ normalized 768-dimensional query embedding via local \texttt{nomic-embed-text} \cite{nussbaum2024nomic}. Performs nearest-neighbor search in FAISS CPU store \cite{johnson2021faiss}, retrieving active/stable memories exceeding threshold $\tau = 0.50$.
    \item \textbf{Node 3: \texttt{sql\_\allowbreak generator\_\allowbreak node}}: Generates initial candidate SQL or regenerates repaired SQL using local \texttt{qwen2.5:7b-instruct} \cite{qwen2024qwen25} under greedy decoding (\texttt{temperature = 0.0}).
    \item \textbf{Node 4: \texttt{ast\_\allowbreak guard\_\allowbreak node}}: Performs static pre-execution Abstract Syntax Tree (AST) validation using SQLGlot \cite{mao2023sqlglot}. Verifies single-statement execution and blocks destructive mutations.
    \item \textbf{Node 5: \texttt{postgres\_\allowbreak executor\_\allowbreak node}}: Submits validated read-only SQL queries to physical PostgreSQL instance, capturing execution status, result rows, execution latency, and raw driver errors.
    \item \textbf{Node 6: \texttt{observation\_\allowbreak node}}: Passively normalizes raw execution exceptions or validation rejections into an immutable \texttt{Runtime\-Observation} record without root-cause classification.
    \item \textbf{Node 7: \texttt{diagnosis\_\allowbreak node}}: Operates deterministically without LLM inference to classify normalized error traces into a canonical 7-tier exception taxonomy, extract broken identifiers, resolve schema candidate remappings, and generate negative constraints.
    \item \textbf{Node 8: \texttt{knowledge\_\allowbreak node}}: Transforms diagnostic findings into an ephemeral \texttt{Runtime\-Knowledge} artifact and evaluates the retry budget ($K \le 3$).
    \item \textbf{Node 9: \texttt{repair\_\allowbreak prompt\_\allowbreak node}}: Assembles the strictly bounded repair prompt containing the isolated \texttt{[STRICT\allowbreak{} REPAIR\allowbreak{} CONSTRAINTS]} block and routes execution back to Node 3.
    \item \textbf{Node 10: \texttt{memory\_\allowbreak governance\_\allowbreak node}}: Applies mathematical governance upon terminal states (\texttt{STATUS\_\allowbreak SUCCESS}, \allowbreak\texttt{STATUS\_\allowbreak FAILED}, \allowbreak\texttt{STATUS\_\allowbreak BLOCKED})[cite: 19]. Enforces mutual exclusion: reinforces applied memories or admits fresh operational knowledge into FAISS.
\end{enumerate}

The complete closed-loop execution topology, deterministic safety barrier, and repair state machine are illustrated in the End-to-End ARMG Architecture and Closed-Loop State Machine diagram.

\begin{figure*}[t]
\centering
\includegraphics[width=\textwidth]{fig1_architecture.png}
\caption{End-to-End ARMG Architecture and Closed-Loop State Machine. The canonical 10-node LangGraph execution graph enforces deterministic AST safety gating prior to PostgreSQL execution, orchestrates bounded 7-tier diagnostic repair ($K \le 3$), and executes algorithmic mutual exclusion between memory reinforcement and admission at terminal outcomes.}
\label{fig:architecture}
\end{figure*}

\section{Runtime Observation and Error Diagnosis}
When an execution error occurs, raw PostgreSQL stderr strings are ingested by the \texttt{RuntimeObserver} and normalized into structured, immutable \texttt{Runtime\-Observation} records.

\subsection{The Canonical 7-Tier Exception Taxonomy}
The \texttt{DiagnosticEngine} deterministically maps normalized observations into a strict 7-tier exception taxonomy, summarized with trigger patterns and candidate repair heuristics in Table~\ref{tab:taxonomy}.

\begin{table*}[t]
\centering
\caption{Canonical 7-Tier Exception Taxonomy Matrix and Candidate Repair Heuristics}
\label{tab:taxonomy}
\footnotesize
\setlength{\tabcolsep}{4.5pt}
\renewcommand{\arraystretch}{1.04}
\begin{tabularx}{\textwidth}{@{}c l l l X@{}}
\toprule
\textbf{Tier} & \textbf{Category} & \textbf{Implementation Definition} & \textbf{Detection Pattern / Trigger} & \textbf{Candidate Repair Action} \\
\midrule
1 & Validation & Syntactic/Safety AST rejections & Non-SELECT, destructive DDL/DML, stacked injections & Abort repair loop; transition to \texttt{STATUS\_BLOCKED} \\
2 & Syntax & Malformed SQL syntax & Driver syntax errors, unclosed quotes, malformed clauses & Strip invalid tokens; regenerate with strict SQL grammar \\
3 & Semantic & Schema/Identifier non-existence & Unknown column/table names, column ambiguity & Remap to closest catalog token; inject negative constraint \\
4 & Planning & Cartesian products / Join failures & Unbounded joins, missing foreign-key predicates & Inject explicit \texttt{JOIN ... ON} clause from schema catalog \\
5 & Permission & Privileged/Administrative commands & Read-only violations, grant/revoke rejections & Block execution; restrict to read-only \texttt{SELECT} \\
6 & Resource & Operational execution timeouts & Query cancellation, memory quota exceeded & Enforce query timeout; suggest predicate pushdown \\
7 & Execution & Unclassified runtime driver failures & Catch-all database exceptions & Fallback to raw normalized driver error trace \\
\bottomrule
\end{tabularx}
\end{table*}

\subsection{Deterministic Candidate Resolution Heuristic}
When a Tier 3 \texttt{Semantic} error occurs on column $c_{\text{broken}}$, candidate replacement identifiers are resolved from active schema tables without LLM inference:
\begin{enumerate}
    \item \textbf{Substring Match}: $+10.0$ bonus if $c_{\text{broken}} \subseteq c_{\text{cand}}$.
    \item \textbf{Token Overlap}: $+8.0 \times |\text{tokens}(c_{\text{broken}}) \cap \text{tokens}(c_{\text{cand}})|$.
    \item \textbf{Sequence Matcher Ratio}: $+4.0 \times \text{difflib.SequenceMatcher.ratio}()$.
    \item \textbf{Data Type \& Domain Affinity}: $+6.0$ bonus if numeric type matches metric intent; $+15.0$ bonus if \texttt{revenue} matches \texttt{gross\_revenue}; $+10.0$ bonus if \texttt{revenue} matches \texttt{net\_profit}.
    \item \textbf{Deterministic Tie-Break}: Highest score descending, then candidate identifier ascending.
\end{enumerate}

\subsection{Ephemeral RuntimeKnowledge Representation}
Diagnostic findings are structured into an immutable \texttt{Runtime\-Knowledge} record containing: \texttt{failure\_\allowbreak type}, \allowbreak\texttt{source\_\allowbreak exception}, \allowbreak\texttt{context}, \allowbreak\texttt{root\_\allowbreak cause}, \allowbreak\texttt{repair\_\allowbreak strategy}, \allowbreak\texttt{negative\_\allowbreak constraints}, \allowbreak\texttt{candidate\_\allowbreak replacements}[cite: 18], initial prior confidence $C_0 = 0.50$, and UTC timestamp.

\section{Memory Governance and Lifecycle}

All governance control equations are implemented in \url{memory/governance.py}. The complete mathematical formulation of all control mechanisms, parameter specifications, and operational roles is summarized in Table~\ref{tab:governance_equations}.

\begin{table*}[t]
\centering
\caption{Mathematical Governance Engine Control Equations and Parameter Specifications}
\label{tab:governance_equations}
\footnotesize
\setlength{\tabcolsep}{4.5pt}
\renewcommand{\arraystretch}{1.04}
\begin{tabularx}{\textwidth}{@{}l l l X@{}}
\toprule
\textbf{Control Mechanism} & \textbf{Formal Equation / Formulation} & \textbf{Parameter Defaults} & \textbf{Operational Implementation Role} \\
\midrule
Operational Utility & $\text{Utility} = C \times \text{SuccessRate} \times \text{ContextSim} \times \text{Recency}$ & Priors: $C_0=0.5, \text{SR}_0=0.5$ & Multi-factor utility evaluation \\
Admission Gating & $\text{Admit}(K) \iff \text{Utility}_0(K) \ge \theta_{\text{admit}}$ & $\theta_{\text{admit}} = 0.25$ & Gating candidate operational knowledge \\
Confidence Escalation & $C_{t+1} = C_t + \alpha (1.0 - C_t)$ & $\alpha = 0.10$ & Asymptotic reinforcement upon success \\
Failure Penalty & $C_{t+1} = \max(0.0, \, C_t \times (1.0 - \beta))$ & $\beta = 0.15$ & Multiplicative confidence penalty on failure \\
Continuous Decay & $C(t) = C_{\text{ref}} \times \exp(-\lambda \Delta t)$ & $\lambda = 0.05\text{ day}^{-1}$ & Exponential decay over time ($*$Unexercised) \\
Mutual Exclusion & $\text{Reinforce}(M) \iff M_{\text{applied}} \neq \emptyset$; $\text{Admit}(K)$ otherwise & Mutually exclusive & Invariant preventing duplicate memory bloat \\
\bottomrule
\multicolumn{4}{@{}p{\linewidth}@{}}{\footnotesize \textsuperscript{*}Temporal decay is fully unit-tested in \url{tests/unit/test_governance.py}, but unexercised under the 3.5-minute benchmark clock.} \\
\end{tabularx}
\end{table*}

\subsection{Multi-Factor Operational Utility}
Operational memory utility is computed via:
\begin{equation}
\text{Utility} = C \times \text{SuccessRate} \times \text{ContextSimilarity} \times \text{Recency}
\end{equation}
where:
\begin{itemize}
    \item $C \in [0.0, 1.0]$ represents confidence score.
    \item $\text{SuccessRate} = \frac{\text{successful\_uses}}{\text{total\_uses}}$ (defaults to prior $0.50$ if unapplied).
    \item $\text{ContextSimilarity} = \frac{1}{1 + d^2} \in [0.0, 1.0]$, derived from FAISS $L_2$ squared distance $d^2$ \cite{johnson2021faiss}.
    \item $\text{Recency} = \frac{1}{1 + \Delta t} \in (0.0, 1.0]$, where $\Delta t \ge 0$ is elapsed time in days/epochs.
\end{itemize}

\subsection{Admission Control}
A newly derived \texttt{Runtime\-Knowledge} instance is admitted to persistent FAISS storage if and only if:
\begin{equation}
\begin{aligned}
\text{Admit}(K) \iff \; & \text{Utility}_0(K) \ge \theta_{\text{admit}} \\
& (\theta_{\text{admit}} = 0.25)
\end{aligned}
\end{equation}
where default candidate prior has $C_0 = 0.50, \text{SuccessRate}_0 = 0.50, \text{ContextSimilarity} = 1.0, \text{Recency} = 1.0 \implies \text{Utility}_0 = 0.25$. Unrecoverable terminal failures exhausting the retry budget are rejected under \texttt{TERMINAL\_\allowbreak FAILURE\_\allowbreak NOT\_\allowbreak ADMITTED}.

\subsection{Asymptotic Confidence Escalation}
When query repair succeeds with an applied memory, confidence escalates asymptotically:
\begin{equation}
C_{t+1} = C_t + \alpha (1.0 - C_t) \quad (\alpha = 0.10)
\end{equation}
transitioning memory from \texttt{NEW} to \texttt{ACTIVE}, and to \texttt{STABLE} when $C_{t+1} \ge 0.80$.

\subsection{Multiplicative Failure Penalty}
If query repair fails after retrieving memory, a penalty is applied:
\begin{equation}
C_{t+1} = \max(0.0, \, C_t \times (1.0 - \beta)) \quad (\beta = 0.15)
\end{equation}
If $C_{t+1} < 0.20$, the memory transitions to \texttt{ARCHIVED}.

\subsection{Continuous Exponential Temporal Decay}
Memory confidence decays continuously over elapsed time:
\begin{equation}
C(t) = C_{\text{ref}} \times \exp(-\lambda \Delta t) \quad (\lambda = 0.05 \text{ day}^{-1})
\end{equation}
The continuous temporal-decay mechanism is mathematically implemented and unit-tested in the software repository; however, because the sequential 25-query benchmark executed in approximately 3.5 minutes ($\Delta t \approx 0.002$ days), temporal decay remained experimentally unexercised. Mode 6 serves as a static no-decay control, and this benchmark does not establish long-horizon decay effectiveness.

\subsection{Algorithmic Mutual Exclusion Invariant}
To eliminate duplicate memory accumulation when a known repair pattern is reused, ARMG enforces:
\begin{equation}
\text{Action} = \begin{cases} 
\text{Reinforce}(M_{\text{applied}}), & \text{if } M_{\text{applied}} \neq \emptyset \\ 
\text{Admit}(K_{\text{new}}), & \text{if } M_{\text{applied}} = \emptyset \\
& \quad \land \; \text{Success} \land K > 0 \\ 
\emptyset, & \text{otherwise} 
\end{cases}
\end{equation}

\subsection{Runtime Memory Lifecycle State Machine}
The lifecycle transitions governing operational memories across their lifespan are depicted in the Runtime Memory Lifecycle State Transition Machine diagram, illustrating admission gating, confidence escalation, failure penalties, archival thresholds, and deletion criteria.

\begin{figure}[t]
\centering
\includegraphics[width=\columnwidth]{fig2_lifecycle.png}
\caption{Runtime Memory Lifecycle State Transition Machine. Operational memories progress across six discrete lifecycle states (\texttt{NEW}, \texttt{ACTIVE}, \texttt{STABLE}, \texttt{DECAYING}, \texttt{ARCHIVED}, \texttt{DELETED}) governed by mathematical admission gating ($\text{Utility}_0 \ge 0.25$), asymptotic confidence escalation ($\alpha = 0.10$), multiplicative penalty ($\beta = 0.15$), and continuous exponential decay ($\lambda = 0.05/\text{day}$).}
\label{fig:lifecycle}
\end{figure}

\section{Runtime-Guided Repair and Negative Constraints}
During repair, the \texttt{RepairSQLGenerator} constructs a structured prompt containing:
\begin{enumerate}
    \item Target database schema context.
    \item Original natural language query.
    \item Previously failed SQL query string.
    \item Normalized PostgreSQL error trace.
    \item \textbf{Strict Negative Constraints}: Explicitly forbidden column names, prohibited join constructs, and banned AST subtrees derived during diagnosis.
    \item \textbf{Operational Memory Context}: Actionable repair rules and root-cause explanations from retrieved memories.
\end{enumerate}

\section{Safety Enforcement}
To enforce pre-execution containment against destructive mutations, ARMG implements a multi-layered guardrail:
\begin{itemize}
    \item \textbf{Prompt Directive}: Instructs the LLM to output only read-only \texttt{SELECT} queries.
    \item \textbf{Deterministic AST Parser}: The \texttt{Execution\-Validator} parses candidate SQL using SQLGlot \cite{mao2023sqlglot} before database driver invocation. Any AST root node matching \texttt{Drop}, \allowbreak\texttt{Delete}, \allowbreak\texttt{Update}, \allowbreak\texttt{Insert}, \allowbreak\texttt{Create}, \allowbreak\texttt{Alter}, \allowbreak\texttt{Truncate\-Table}, \allowbreak\texttt{Command}, \allowbreak\texttt{Transaction}, \allowbreak\texttt{Commit}, or \allowbreak\texttt{Rollback} is immediately rejected. Administrative keywords (\texttt{GRANT}, \allowbreak\texttt{REVOKE}, \allowbreak\texttt{MERGE}, \allowbreak\texttt{EXEC}) and multi-statement injections (statement count $> 1$) are strictly forbidden.
    \item \textbf{Blocked State Containment}: Safety rejections transition the state graph directly to \texttt{STATUS\_\allowbreak BLOCKED}, bypassing PostgreSQL execution entirely and rejecting memory admission.
\end{itemize}

\section{Experimental Methodology}

\subsection{Evaluation Configuration}
\begin{itemize}
    \item \textbf{Model}: \texttt{qwen2.5:7b-instruct} \cite{qwen2024qwen25} (Alibaba Cloud / Qwen, 7.61B parameters, local Ollama instance). Adaptation is strictly inference-time based without parameter updates or fine-tuning.
    \item \textbf{Decoding Configuration}: Greedy decoding (\texttt{temperature = 0.0}) enforced across initial generation and repair generation.
    \item \textbf{Embedding Model}: \texttt{nomic-embed-text} \cite{nussbaum2024nomic} (768 dimensions, Unit-$L_2$ normalized).
    \item \textbf{Vector Index}: FAISS \texttt{IndexIDMap2} wrapping \texttt{IndexFlatL2(768)} \cite{johnson2021faiss}.
    \item \textbf{Database}: PostgreSQL 16 on \texttt{localhost:5432} (\texttt{armg\_db}).
    \item \textbf{Benchmark Warehouse}: A synthetically generated enterprise-style Star Schema Data Warehouse modeling B2B technology product transactions across calendar year 2025, populated deterministically via NumPy \texttt{seed=42} (\url{scripts/seed_warehouse.py}, 2,000 fact records, 4 tables: \texttt{dim\_time}, \allowbreak\texttt{dim\_geography}, \allowbreak\texttt{dim\_product}, \allowbreak\texttt{fact\_\allowbreak sales\_\allowbreak performance}). We utilize a controlled warehouse to enable precise measurement of multi-turn operational memory and lifecycle governance, contrasting with cross-domain academic benchmarks (Spider \cite{yu2018spider}, BIRD \cite{li2023bird}) that evaluate single-turn schema generalizability across hundreds of independent databases.
    \item \textbf{Benchmark Corpus}: 25 analytical queries (\url{benchmark/queries.json}) across Categories A, B, C, D.
    \item \textbf{Repeated Experimental Protocol}: Three isolated repeated benchmark executions (Seeds 42, 123, 999) across six experimental modes ($3 \times 6 \times 25 = 450$ total evaluated query runs). Because greedy decoding was enforced and benchmark seeds were not passed to Ollama, these repeated runs evaluate \textbf{pipeline reproducibility, execution stability, and memory-state consistency}, rather than stochastic sampling variance from 450 independent random trials ($n = 3$ repeated runs).
\end{itemize}

The canonical system configuration, component mapping, and operational specifications are summarized in Table~\ref{tab:system_config}. The architectural specification, table roles, row counts, primary keys, and foreign-key integrity constraints for the synthetic Star Schema warehouse are detailed in Table~\ref{tab:schema}. The complexity distribution, analytical focus, and query clause characteristics across the 25 benchmark queries are presented in Table~\ref{tab:query_corpus}.

\begin{table}[t]
\centering
\caption{Canonical System Configuration and Component Mapping}
\label{tab:system_config}
\footnotesize
\setlength{\tabcolsep}{2.5pt}
\renewcommand{\arraystretch}{1.04}
\begin{tabularx}{\columnwidth}{@{}p{2.1cm} p{2.1cm} X@{}}
\toprule
\textbf{Subsystem} & \textbf{Anchor} & \textbf{Specification} \\
\midrule
Foundation Model & Local Ollama & \texttt{qwen2.5:7b-instruct} (7.61B) \\
Inference Config & Local Ollama & Greedy (\texttt{temp=0.0}, \texttt{top\_p=1.0}) \\
Embedding Model & Local Ollama & \texttt{nomic-embed-text} (768d, Unit-$L_2$) \\
Vector Index & CPU Flat Index & FAISS \texttt{IndexIDMap2} (\texttt{IndexFlatL2}) \\
Warehouse & Container & PostgreSQL 16.2 on \texttt{localhost:5432} \\
Graph Engine & Directed Graph & LangGraph 0.2.x (\texttt{StateGraph}) \\
Static Parser & AST Guardrail & SQLGlot 25.x (Dialect: PostgreSQL) \\
Driver & DB-API 2.0 & \texttt{psycopg2-binary} 2.9.x \\
Runtime & Workstation & Python 3.11.9 (CUDA via Ollama) \\
\bottomrule
\end{tabularx}
\end{table}

\begin{table*}[t]
\centering
\caption{Relational Data Warehouse Star Schema Specification}
\label{tab:schema}
\footnotesize
\setlength{\tabcolsep}{5pt}
\renewcommand{\arraystretch}{1.04}
\begin{tabularx}{\textwidth}{@{}l c c l X@{}}
\toprule
\textbf{Table Name} & \textbf{Role} & \textbf{Rows} & \textbf{Primary Key} & \textbf{Major Attributes / Foreign Key Constraints} \\
\midrule
\texttt{dim\_time} & Dimension & 365 & \texttt{time\_key} & full\_date, day\_of\_week, calendar\_month, calendar\_quarter, calendar\_year \\
\texttt{dim\_geography} & Dimension & 6 & \texttt{geo\_key} & region, zone, market\_type \\
\texttt{dim\_product} & Dimension & 8 & \texttt{product\_key} & product\_name, category, sub\_category, unit\_cost \\
\texttt{fact\_sales\_performance} & Fact & 2,000 & \texttt{fact\_key} & units\_sold, gross\_revenue, discount\_applied, net\_profit. \newline FK: \texttt{time\_key}, \texttt{geo\_key}, \texttt{product\_key} \\
\bottomrule
\multicolumn{5}{@{}p{\linewidth}@{}}{\footnotesize \textsuperscript{*}Deterministically seeded via NumPy \texttt{seed=42} (\url{scripts/seed_warehouse.py}).} \\
\end{tabularx}
\end{table*}

\begin{table}[t]
\centering
\caption{Benchmark Query Corpus Distribution Across Complexity Categories}
\label{tab:query_corpus}
\footnotesize
\setlength{\tabcolsep}{3pt}
\renewcommand{\arraystretch}{1.04}
\begin{tabularx}{\columnwidth}{@{}l c c X@{}}
\toprule
\textbf{Category} & \textbf{Queries} & \textbf{Count} & \textbf{Analytical Focus and SQL Clause Complexity} \\
\midrule
Category A & Q01--Q05 & 5 & Simple aggregations, basic filters, group-by, order-by clauses \\
Category B & Q06--Q13 & 8 & Multi-table Star Schema joins, dimension filtering, compound conditions \\
Category C & Q14--Q19 & 6 & Advanced window functions (\texttt{RANK()}, \texttt{LAG()}, cumulative partitions) \\
Category D & Q20--Q25 & 6 & Semantic/schema trap queries, attribute sequence inversions, strict ordering \\
\midrule
\textbf{Total Corpus} & Q01--Q25 & 25 & B2B Technology Sales Analytics Domain (\url{benchmark/queries.json}) \\
\bottomrule
\end{tabularx}
\end{table}

\subsection{The Six Experimental Modes}
\begin{enumerate}
    \item \textbf{Mode 1 (Zero-Shot)}: Monolithic single-pass prompt; zero repair; zero memory ($K=0$).
    \item \textbf{Mode 2 (Stateless Self-Correction)}: Iterative repair ($K \le 3$) feeding raw error strings; memoryless.
    \item \textbf{Mode 3 (Naive Vector RAG)}: Appends raw \texttt{(question, sql)} pairs upon success; retrieves top-3 few-shot examples; zero repair; zero governance.
    \item \textbf{Mode 4 (Full ARMG)}: Governed retrieval, 7-tier diagnosis, negative constraints, bounded repair ($K \le 3$), mutual-exclusion governance.
    \item \textbf{Mode 5 (ARMG $-$ Negative Constraints)}: Identical to Mode 4, but \texttt{[STRICT\allowbreak{} REPAIR\allowbreak{} CONSTRAINTS]} block is omitted.
    \item \textbf{Mode 6 (ARMG with $\lambda = 0.0$)}: Identical to Mode 4, but continuous decay rate is set to 0.0 (static no-decay control).
\end{enumerate}

\subsection{Relational Equivalence Comparator Rules}
Following best practices in Text-to-SQL evaluation methodology \cite{finegandollak2018eval, zhong2020distilled}, the relational equivalence comparator enforces: (1) Execution Failure Gate; (2) Cardinality and Dimensionality equality; (3) Empty Set Handling; (4) Ordering Detection; and (5) Numerical Precision tolerance ($\vert{}x - y\vert{} \le 10^{-4}$).

\section{Experimental Results}

\subsection{Comparative Empirical Results across Six Modes}
Table~\ref{tab:results} presents descriptive empirical results averaged across the three repeated benchmark executions ($n = 3$ repeated runs).

\begin{table*}[t]
\centering
\caption{Comparative Empirical Benchmark Results Across Six Experimental Modes ($n = 3$ Repeated Executions)}
\label{tab:results}
\footnotesize
\setlength{\tabcolsep}{8pt}
\renewcommand{\arraystretch}{1.05}
\begin{tabular}{lcccccc}
\toprule
\textbf{Experimental Mode} & \textbf{ExecAcc (\%)} & \textbf{ExecSucc (\%)} & \textbf{Retries ($K$)} & \textbf{Latency (ms)} & \textbf{Tokens} & \textbf{Store Size} \\
\midrule
Mode 1 (Zero-Shot) & $57.33 \pm 2.31$ & $76.00 \pm 0.00$ & $0.00 \pm 0.00$ & $5,087.73 \pm 210.33$ & $360.48 \pm 0.48$ & $0$ \\
Mode 2 (Stateless Self-Correction) & $68.00 \pm 0.00$ & $92.00 \pm 0.00$ & $0.43 \pm 0.02$ & $7,149.68 \pm 231.65$ & $572.72 \pm 12.68$ & $0$ \\
Mode 3 (Naive Vector RAG) & $68.00 \pm 0.00$ & $92.00 \pm 0.00$ & $0.00 \pm 0.00$ & $6,844.01 \pm 56.90$ & $558.28 \pm 0.00$ & $23$ \\
Mode 4 (Full ARMG) & $68.00 \pm 0.00$ & $96.00 \pm 0.00$ & $0.28 \pm 0.00$ & $8,564.89 \pm 103.16$ & $542.37 \pm 0.02$ & $3$ \\
Mode 5 (ARMG $-$ Neg Constraints) & $68.00 \pm 0.00$ & $96.00 \pm 0.00$ & $0.36 \pm 0.00$ & $8,941.28 \pm 108.68$ & $553.93 \pm 0.02$ & $3$ \\
Mode 6 (ARMG with $\lambda = 0.0$) & $68.00 \pm 0.00$ & $94.67 \pm 2.31$ & $0.37 \pm 0.02$ & $9,036.97 \pm 178.55$ & $599.67 \pm 13.94$ & $3$ \\
\bottomrule
\multicolumn{7}{@{}p{\linewidth}@{}}{\footnotesize \textsuperscript{*}All reported figures represent mean $\pm$ sample standard deviation across three repeated executions under seeds 42, 123, and 999.} \\
\end{tabular}
\end{table*}

The PostgreSQL Execution Success vs. Relational Execution Accuracy comparison illustrates the distribution across all six evaluated experimental modes, demonstrating the consistent divergence between syntactically executable queries and semantic relational accuracy.

\begin{figure}[t]
\centering
\includegraphics[width=\columnwidth]{fig4_execsucc_execacc.png}
\caption{PostgreSQL Execution Success vs. Relational Execution Accuracy across six experimental modes ($n = 3$ repeated runs). Shaded bars illustrate PostgreSQL execution success ($\text{ExecSucc}$), while dark hatched bars represent relational semantic accuracy ($\text{ExecAcc}$). Across Modes 2--6, relational execution accuracy plateaus at 68.00\% despite execution success reaching up to 96.00\%, illustrating the critical gap between execution and semantic equivalence.}
\label{fig:exec_gap}
\end{figure}

\subsection{Mode 4 vs. Mode 2 Performance Trade-Off Analysis}
The multi-dimensional trade-off profile between Mode 2 (Stateless Self-Correction) and Mode 4 (Full ARMG) is quantitatively summarized in Table~\ref{tab:tradeoff}.

\begin{table*}[t]
\centering
\caption{Mode 4 (Full ARMG) vs. Mode 2 (Stateless Self-Correction) Comparative Trade-Off Profile}
\label{tab:tradeoff}
\footnotesize
\setlength{\tabcolsep}{4.5pt}
\renewcommand{\arraystretch}{1.04}
\begin{tabularx}{\textwidth}{@{}l r r c c X@{}}
\toprule
\textbf{Evaluation Dimension} & \textbf{Mode 2} & \textbf{Mode 4} & \textbf{Absolute Delta} & \textbf{Relative Delta} & \textbf{Operational Engineering Interpretation} \\
\midrule
Mean Repair Retries ($K$) & 0.43 & 0.28 & $-0.15$ retries & $-34.38\%$ & Observed reduction in repair iterations \\
Mean Token Expenditure & 572.72 & 542.37 & $-30.35$ tokens & $-5.30\%$ & Token savings by avoiding repetitive repair prompts \\
PostgreSQL Execution Success & $92.00\%$ & $96.00\%$ & $+4.00\text{ pp}$ & $+4.35\%$ & Execution recovery on Query Q19 \\
End-to-End Latency & $7,149.68\text{ ms}$ & $8,564.89\text{ ms}$ & $+1,415.21\text{ ms}$ & $+19.79\%$ & Architectural overhead of state graph and FAISS \\
Relational Semantic Accuracy & $68.00\%$ & $68.00\%$ & $0.00\text{ pp}$ & $0.00\%$ & Observed identical relational execution accuracy in the evaluated benchmark \\
\bottomrule
\multicolumn{6}{@{}p{\linewidth}@{}}{\footnotesize \textsuperscript{*}Percentage points (pp) and relative percentage changes (\%) are strictly distinguished.} \\
\end{tabularx}
\end{table*}

The Mode 2 vs. Full ARMG Operational Trade-Off Profile visualizes this multi-dimensional comparison across the five core operational metrics, explicitly distinguishing relative percentage changes from percentage-point shifts.

\begin{figure}[t]
\centering
\includegraphics[width=\columnwidth]{fig5_tradeoff.png}
\caption{Mode 2 vs. Full ARMG Operational Trade-Off Profile. Demonstrates the observed operational trade-offs of Full ARMG relative to stateless self-correction: a 34.38\% reduction in repair iterations, 5.30\% token reduction, and +4.00 percentage points in execution success, balanced against a 19.79\% latency overhead, with relational accuracy remaining identical (0.00 pp).}
\label{fig:tradeoff}
\end{figure}

\begin{itemize}
    \item \textbf{Repair Iteration Efficiency}: Mode 4 achieved a \textbf{34.38\% reduction in mean repair iterations} (0.28 vs. 0.43 retries per query), observed consistently across all three runs (Seed 42: 0.28 vs. 0.44; Seed 123: 0.28 vs. 0.44; Seed 999: 0.28 vs. 0.40).
    \item \textbf{Token Expenditure}: Mode 4 consumed \textbf{5.30\% fewer tokens} (542.37 vs. 572.72 tokens per query) by avoiding repetitive repair prompt/completion cycles.
    \item \textbf{PostgreSQL Execution Success}: Mode 4 achieved \textbf{96.00\% execution success} vs. 92.00\% for Mode 2 (+4.00 percentage points; 24/25 vs. 23/25 queries).
    \item \textbf{Latency Overhead}: In the evaluated benchmark, Mode 4 exhibited a \textbf{19.79\% higher end-to-end latency} (8,564.89~ms vs. 7,149.68~ms) while using fewer mean repair iterations, which coincided with the additional embedding, retrieval, and state-graph orchestration stages.
    \item \textbf{Relational Accuracy Plateau}: Both Mode 4 and Mode 2 achieved exactly \textbf{68.00\% relational execution accuracy} (17 of 25 queries correct).
\end{itemize}

\subsection{Query-Level Divergence Analysis}
Across all 25 queries, Mode 4 and Mode 2 differed on exactly two queries:
\begin{enumerate}
    \item \textbf{Query Q08 (Category B --- Join Aggregation)}: Mode 2 failed on Attempt 1 and required a repair retry (mean 0.67 retries across runs). Mode 4 retrieved memory admitted during Q04 and executed cleanly on Attempt 1 without retries (\texttt{retries = 0}).
    \item \textbf{Query Q19 (Category C --- Percentage Contribution)}: Mode 2 encountered a \texttt{GROUP BY} syntax error, exhausted all 3 retries, and failed (\texttt{retries = 3}, \texttt{is\_success = False}). Mode 4 retrieved memory admitted during Q13 and generated executable SQL on Attempt 1 (\texttt{retries = 0}, \texttt{is\_success = True}).
    \item \textbf{Remaining 23 Queries}: Exhibited identical execution outcomes across both modes.
\end{enumerate}

\section{Memory Retrieval and Lifecycle Analysis}

\subsection{Remediation of Vector Geometry}
In historical pre-remediation testing, raw unnormalized embeddings generated by \texttt{nomic-embed-text} had norms $\|\mathbf{v}\| \approx 19.8$, resulting in squared $L_2$ distances $d^2 \approx 280$ and similarity scores $S = \frac{1}{1 + d^2} \approx 0.0035 \ll 0.50$. Consequently, zero retrievals occurred across all queries. Enforcing Unit-$L_2$ normalization restored the intended similarity geometry, yielding 16 retrieval events across 12 distinct queries ($48.0\%$ benchmark coverage).

The Pre-Remediation vs. Post-Remediation FAISS Retrieval Geometry distribution illustrates similarity scores before and after embedding normalization, demonstrating the restoration of the intended geometric threshold $\tau = 0.50$.

\begin{figure*}[t]
\centering
\includegraphics[width=\textwidth]{fig3_retrieval_geometry.png}
\caption{Pre-Remediation vs. Post-Remediation FAISS Retrieval Geometry. Left: Unnormalized \texttt{nomic-embed-text} embeddings generated large norms ($\Vert{}\mathbf{v}\Vert{} \approx 19.8$), driving all similarity scores to $S \approx 0.0035 \ll \tau = 0.50$ (zero retrievals). Right: Unit-$L_2$ normalization restored proper inner-product geometry, elevating 16 retrieval events above threshold $\tau = 0.50$ across 12 distinct queries ($48.0\%$ coverage).}
\label{fig:retrieval_geometry}
\end{figure*}

\subsection{Lifecycle Metrics and Deduplication}
\begin{itemize}
    \item \textbf{New Admissions}: Exactly 3 memories were admitted across all runs: \texttt{mem-Q04} (admitted join pattern for \texttt{dim\_geography}), \texttt{mem-Q13} (admitted multi-table join and aggregation pattern), and \texttt{mem-Q15} (admitted window aggregation structure).
    \item \textbf{Reinforcement \& Deduplication on \texttt{Q17}}: Q17 applied and reinforced \texttt{mem-Q13}; new admission was suppressed by the mutual-exclusion rule. Specifically, for Q17, the repaired query succeeded on Attempt 2 after retrieved operational memory was applied, and the applied memory was subsequently reinforced.
    \item \textbf{Store Stability}: The evaluated Mode 4 store remained at 3 memories from Q15 through Q25 across all three evaluated runs, whereas the evaluated Naive Vector RAG control accumulated 23 entries.
\end{itemize}

The Persistent Memory Store Growth trajectory illustrates cumulative memory accumulation across the sequential benchmark queries, comparing Full ARMG against Naive Vector RAG.

\begin{figure}[t]
\centering
\includegraphics[width=\columnwidth]{fig6_memory_growth.png}
\caption{Persistent Memory Store Growth: Full ARMG vs. Naive Vector RAG. In Full ARMG (Mode 4), algorithmic mutual exclusion between memory reinforcement and admission maintained the persistent store size strictly invariant at exactly 3 memories from Q15 through Q25. In contrast, unmanaged Naive Vector RAG (Mode 3) appended uncurated exemplars upon every execution success, expanding to 23 memories.}
\label{fig:memory_growth}
\end{figure}

\section{Safety Evaluation}
Across 450 post-remediation evaluations and 150 historical baseline evaluations (600 audited executions in total), no destructive SQL statement reached PostgreSQL in the evaluated benchmark.

In offline safety verification testing (\url{tests/unit/test_safety_guard.py}, 18 unit tests), when presented with adversarial requests (\texttt{DROP TABLE}, \allowbreak\texttt{DELETE FROM}, \allowbreak\texttt{UPDATE}, stacked injections, administrative grants), the AST guard halted all tested destructive-query cases before database submission, classifying the violations as \texttt{destructive\_mutation}, terminating the repair loop, and transitioning directly to \texttt{STATUS\_\allowbreak BLOCKED}.

\section{Semantic Failure Analysis}

A central methodological finding is the substantial divergence between PostgreSQL execution success and relational execution accuracy across all six experimental modes, as detailed in Table~\ref{tab:gap}.

\begin{table}[t]
\centering
\caption{PostgreSQL Execution Success vs. Relational Semantic Accuracy Across All Six Modes}
\label{tab:gap}
\footnotesize
\setlength{\tabcolsep}{6pt}
\renewcommand{\arraystretch}{1.06}
\begin{tabular}{lccc}
\toprule
\textbf{Experimental Mode} & \textbf{ExecSucc (\%)} & \textbf{ExecAcc (\%)} & \textbf{Gap (pp)} \\
\midrule
Mode 1 (Zero-Shot) & 76.00 & 57.33 & 18.67 \\
Mode 2 (Stateless Self-Correction) & 92.00 & 68.00 & 24.00 \\
Mode 3 (Naive Vector RAG) & 92.00 & 68.00 & 24.00 \\
Mode 4 (Full ARMG) & 96.00 & 68.00 & 28.00 \\
Mode 5 (ARMG $-$ Neg Constraints) & 96.00 & 68.00 & 28.00 \\
Mode 6 (ARMG with $\lambda = 0.0$) & 94.67 & 68.00 & 26.67 \\
\bottomrule
\multicolumn{4}{@{}p{\linewidth}@{}}{\scriptsize Discrepancy Gap defined as $\text{ExecSucc} - \text{ExecAcc}$ in percentage points (pp).} \\
\end{tabular}
\end{table}

\subsection{Analysis of the 7 Divergent Queries in Mode 4}
Seven queries were executable but divergent under the relational comparator; four involved advanced window-function constructs, while the remaining divergences involved ordering or projection requirements. A detailed forensic breakdown of these seven divergent queries in Mode 4 under the frozen relational comparator is presented in Table~\ref{tab:divergent_queries}, cataloging the specific SQL construct deviations and comparator diagnostic rationales.

\begin{table*}[t]
\centering
\caption{Forensic Diagnostic Breakdown of the Seven Divergent Queries in Mode 4 Under the Frozen Relational Comparator}
\label{tab:divergent_queries}
\footnotesize
\setlength{\tabcolsep}{4pt}
\renewcommand{\arraystretch}{1.05}
\begin{tabularx}{\textwidth}{@{}c c c c p{4.2cm} X@{}}
\toprule
\textbf{Query} & \textbf{Cat.} & \textbf{PG Status} & \textbf{Rel. Acc.} &
\textbf{Generated SQL Construct Deviation} &
\textbf{Comparator Diagnostic Rationale} \\
\midrule
Q05$^{*}$ & A & Success & False &
Omitted required \texttt{ORDER BY net\_profit DESC} &
Comparator required the same ordered tuple sequence as the gold query; PostgreSQL execution itself succeeded \\
Q14 & C & Success & False &
Substituted simple \texttt{ORDER BY} for \texttt{RANK() OVER} &
Failed multiset row-ranking equivalence \\
Q15 & C & Success & False &
Computed monthly aggregation without a cumulative frame &
Omitted the required running-total window specification \\
Q17 & C & Success & False &
Included an invalid grouping attribute in \texttt{LAG()} partition &
Produced multi-row monthly output instead of the expected scalar lag result \\
Q18 & C & Success & False &
Ranked globally without \texttt{PARTITION BY category} &
Missed the required category-scoped partitioning \\
Q19 & C & Success & False &
Omitted the base revenue column and required rounding format &
Projection signature and decimal-precision discrepancy \\
Q25 & D & Success & False &
Inverted column sequence: \texttt{(market, profit, rev)} &
Positional tuple-attribute mismatch against the gold signature \\
\bottomrule
\multicolumn{6}{@{}p{\linewidth}@{}}{\footnotesize
\textsuperscript{*}Q05 executed successfully on PostgreSQL but was classified as relationally inaccurate because the frozen comparator enforces strict positional sequence matching when the gold query specifies ordering.
} \\
\end{tabularx}
\end{table*}

\begin{enumerate}
    \item \texttt{Q05} (Category A): Q05 executed successfully on PostgreSQL but was classified as relationally inaccurate by the frozen comparator because the generated query omitted the ordering specified by the gold query (\texttt{ORDER BY net\_profit DESC}). When gold SQL specifies ordering, the comparator enforces positional tuple sequence matching, which the unordered result set failed, despite clean database execution.
    \item \texttt{Q14} (Category C): Replaced the \texttt{RANK() OVER (ORDER BY revenue DESC)} window function with a simple \texttt{ORDER BY} clause.
    \item \texttt{Q15} (Category C): Omitted cumulative window framing, computing simple monthly aggregations rather than a running total.
    \item \texttt{Q17} (Category C): Included an invalid grouping attribute in the \texttt{LAG()} partition, generating multi-row monthly output.
    \item \texttt{Q18} (Category C): Omitted \texttt{PARTITION BY category} in the \texttt{RANK()} function, ranking globally across the entire table.
    \item \texttt{Q19} (Category C): Generated executable SQL but computed the percentage contribution without projecting the required base revenue column and \texttt{ROUND(..., 2)} formatting.
    \item \texttt{Q25} (Category D): Projected columns in inverted sequence \texttt{(market\_type, profit, revenue)} instead of \texttt{(market\_type, revenue, profit)}.
\end{enumerate}

\subsection{Root Cause Interpretation \& Model Scale Boundaries}
In the evaluated benchmark, ARMG did not improve relational execution accuracy beyond the 68.00\% level observed for stateless self-correction, including on the evaluated window-function queries. In the evaluated Qwen2.5 7B configuration, remaining failures were concentrated in complex analytical window constructs (4 queries) and projection/ordering specifications (3 queries). While existing model-scaling literature provides broader context on capability variation with parameter scale \cite{wei2022emergent}, the present benchmark does not causally isolate model size. Rather, the empirical results demonstrate that within the evaluated 7B setting, operational memory provides syntactic and error-avoidance guidance, but in the evaluated benchmark, operational memory did not overcome the observed semantic errors on the more complex window-function queries in the local 7B configuration.

\section{Ablation Analysis}

\subsection{Negative Constraints Ablation (Mode 4 vs. Mode 5)}
Removing negative constraints increased mean retries from 0.28 to 0.36 (+28.57\%). Query-level inspection reveals that this entire difference was localized to \textbf{Query \texttt{Q19}}:
\begin{itemize}
    \item In Mode 4 (with negative constraints), the model avoided invalid grouping constructs and executed on Attempt 1 (\texttt{retries = 0}).
    \item In Mode 5 (without negative constraints), Attempt 1 repeated an invalid grouping construct, requiring 2 repair retries (\texttt{retries = 2}).
    \item \textit{Finding}: The negative-constraint ablation showed an observed retry difference localized to Q19; the benchmark does not establish a general effect across queries.
\end{itemize}

\subsection{Temporal Decay Ablation (Mode 4 vs. Mode 6)}
Mode 6 ($\lambda = 0.0$) exhibited identical final store sizes (3 memories) and retrieval counts (16 retrievals) to Mode 4. Because the 25 benchmark queries execute in continuous sequence within $\sim$3.5 minutes ($\Delta t \approx 0.002$ days), exponential decay factor $\exp(-\lambda \Delta t) \approx 0.9999$ was insufficient to age memories.
\begin{itemize}
    \item \textit{Finding}: \textbf{Temporal decay effectiveness was NOT demonstrated by this benchmark}. Mode 6 functioned as a static no-decay control, and this benchmark does not establish long-horizon decay effectiveness.
\end{itemize}

\section{Discussion}

\subsection{What Improved}
\begin{itemize}
    \item \textbf{Repair Efficiency}: ARMG reduced mean repair retries by 34.38\% (0.28 vs. 0.43 retries per query) in the evaluated benchmark, avoiding repetitive repair cycling on recurring schema traps.
    \item \textbf{Token Economy}: Bounded repair behavior achieved a 5.30\% reduction in mean token expenditure (542.37 vs. 572.72 tokens).
    \item \textbf{Execution Success}: PostgreSQL execution success improved from 92.00\% to 96.00\%, recovering an executable query on Q19.
    \item \textbf{Memory Lifecycle Control}: In the evaluated sequential benchmark, algorithmic mutual exclusion between reinforcement and new admission maintained persistent store size at 3 memories after Q15, whereas the evaluated Naive Vector RAG control accumulated 23 entries.
    \item \textbf{Verified Execution Safety}: The AST guard halted all tested destructive-query cases before database submission across unit and benchmark evaluations; no destructive SQL statements reached PostgreSQL in the evaluated benchmark.
\end{itemize}

\subsection{What Did Not Improve}
\begin{itemize}
    \item \textbf{Relational Execution Accuracy}: In the evaluated benchmark, ARMG did not improve relational execution accuracy beyond the 68.00\% level observed for stateless self-correction (17 of 25 queries). Operational memory provided operational and syntactic guidance, but did not overcome the observed semantic errors on the evaluated window-function queries in the local 7B configuration.
\end{itemize}

\subsection{Architectural Cost: Latency Overhead}
\begin{itemize}
    \item In the evaluated benchmark, ARMG exhibited a 19.79\% higher end-to-end latency (8,564.89~ms vs. 7,149.68~ms) while using fewer mean repair iterations. The measured latency difference coincided with the additional embedding, retrieval, and state-graph processing stages.
\end{itemize}

\subsection{What Remains Untested}
\begin{itemize}
    \item \textbf{Long-Term Temporal Decay}: The continuous temporal-decay mechanism remains experimentally unexercised over multi-week or multi-month operational epochs under live calendar drift.
    \item \textbf{Causal Memory Attribution}: While observational associations were verified on Q08 and Q19, establishing formal causal proof requires counterfactual memory-masking interventions.
    \item \textbf{Broader Generalization}: Performance across larger models (70B+) and public multi-schema benchmarks (Spider, BIRD) remains to be established.
\end{itemize}

\section{Limitations}
\begin{enumerate}
    \item \textbf{Model Parameter Scale}: Evaluated exclusively using a local 7B open-weights model (\texttt{qwen2.5:7b-instruct}). Frontier models may exhibit different baseline repair dynamics.
    \item \textbf{Benchmark Corpus Scale}: Evaluated over a fixed corpus of 25 analytical queries over a 4-table Star Schema warehouse.
    \item \textbf{Data Scope}: Evaluated over a synthetically seeded warehouse (2,000 fact records) rather than live production enterprise data.
    \item \textbf{Deterministic Decoding}: Greedy decoding (\texttt{temperature = 0.0}) evaluates deterministic pipeline stability rather than stochastic sampling distributions.
    \item \textbf{Replication Sample Size}: $n = 3$ repeated benchmark executions over 25 queries, evaluating pipeline stability rather than stochastic sampling variance from 450 independent random trials.
    \item \textbf{Temporal Decay Scope}: Rapid benchmark execution clock ($\sim$3.5 minutes) did not exercise continuous exponential decay.
    \item \textbf{Cross-Domain Benchmarks}: Not evaluated on public cross-domain benchmarks (Spider, BIRD).
    \item \textbf{Absence of Counterfactual Intervention}: Dynamic memory masking was not performed during inference.
    \item \textbf{Workload Model}: Evaluated under a single-user sequential query workload without concurrent multi-user load.
\end{enumerate}

\section{Threats to Validity}
\begin{itemize}
    \item \textbf{Internal Validity}: Greedy decoding ensures execution reproducibility, but precludes evaluating temperature-dependent variance. Procedural query sequencing ($Q01 \to Q25$) allows memory transfer from earlier to later queries, modeling realistic operational sessions but introducing order dependency.
    \item \textbf{Construct Validity}: PostgreSQL execution success does not imply relational semantic equivalence. The relational equivalence comparator mitigates this by enforcing multiset bag equivalence and strict positional matching when \texttt{ORDER BY} is required, but does not perform symbolic AST proof.
    \item \textbf{External Validity}: Results are established on a single Star Schema data warehouse. Generalization to enterprise schemas with hundreds of normalized tables or non-relational datastores remains unverified.
    \item \textbf{Statistical Validity}: With $n = 3$ repeated runs over 25 queries, inferential tests ($t$-tests, ANOVAs) are underpowered. All reported results represent descriptive empirical effect sizes across repeated executions.
    \item \textbf{Reproducibility}: High. All random seeds, queries, and execution parameters are fully locked in the repository.
\end{itemize}

\section{Future Work}
\begin{enumerate}
    \item \textbf{Synthetic Epoch Advance Decay Experiment}: Augment evaluation with simulated multi-month time jumps to evaluate continuous exponential decay and archival pruning.
    \item \textbf{Counterfactual Memory Intervention}: Execute controlled A/B testing dynamically masking retrieved memories for Q08 and Q19 to establish formal causal proof of memory-driven repair.
    \item \textbf{Frontier Model Evaluation}: Evaluate ARMG with larger models (Llama-3-70B, Qwen-2.5-72B) to assess whether higher reasoning capacity breaks the 68.00\% semantic window-function ceiling.
    \item \textbf{Academic Benchmark Evaluation}: Port ARMG to Spider and BIRD benchmarks to establish comparative performance against published literature leaders.
    \item \textbf{Concurrent Throughput Testing}: Benchmark FAISS retrieval and database connection pooling under concurrent multi-user load.
\end{enumerate}

\section{Conclusion}
This paper presented Adaptive Runtime Memory Governance (ARMG), an operational knowledge framework for local Text-to-SQL systems. Across three repeated executions of a 25-query benchmark, ARMG demonstrated verified retrieval restoration, mutual-exclusion duplicate suppression, pre-execution safety containment preventing destructive statements from reaching PostgreSQL, an observed 34.38\% reduction in repair loop iterations, and a 5.30\% reduction in token consumption compared to stateless self-correction. Simultaneously, the evaluation established that ARMG incurred a 19.79\% latency overhead and did not improve relational execution accuracy beyond the 68.00\% level observed for stateless self-correction on the evaluated queries. ARMG provides a principled, governed operational memory architecture for enterprise Text-to-SQL systems prioritizing pre-execution safety and repair iteration efficiency within evaluated local settings.

\appendix
\section{Numerical Validation Tables Underlying Primary Diagrams}

This appendix provides the full numerical tables underlying the primary manuscript diagrams, ensuring complete empirical traceability against the raw repository execution logs.

\begin{table*}[t]
\centering
\caption{Numerical Data Underlying the FAISS Retrieval Geometry Analysis: Pre- and Post-Remediation Similarity Scores and Retrieval Events for the 25-Query Benchmark}
\label{tab:fig3_validation}
\footnotesize
\setlength{\tabcolsep}{16pt}
\renewcommand{\arraystretch}{1.03}
\begin{tabular}{l c c c c}
\toprule
\textbf{Query} & \textbf{Pre-remediation $S$} & \textbf{Post-remediation $S$} & \textbf{Pre Retrieval} & \textbf{Post Retrieval Events} \\
\midrule
Q01 & 0.0036 & 0.3608 & No (0) & No (0) \\
Q02 & 0.0035 & 0.3171 & No (0) & No (0) \\
Q03 & 0.0037 & 0.3065 & No (0) & No (0) \\
Q04 & 0.0040 & 0.3949 & No (0) & No (0) \\
Q05 & 0.0034 & 0.3966 & No (0) & No (0) \\
Q06 & 0.0034 & 0.3808 & No (0) & No (0) \\
Q07 & 0.0040 & 0.3305 & No (0) & No (0) \\
Q08 & 0.0037 & 0.6459 & No (0) & Yes (1) \\
Q09 & 0.0034 & 0.3684 & No (0) & No (0) \\
Q10 & 0.0037 & 0.6664 & No (0) & Yes (1) \\
Q11 & 0.0034 & 0.3122 & No (0) & No (0) \\
Q12 & 0.0034 & 0.3495 & No (0) & No (0) \\
Q13 & 0.0036 & 0.3034 & No (0) & No (0) \\
Q14 & 0.0029 & 0.3909 & No (0) & No (0) \\
Q15 & 0.0030 & 0.3259 & No (0) & No (0) \\
Q16 & 0.0033 & 0.6798 & No (0) & Yes (1) \\
Q17 & 0.0032 & 0.8125 & No (0) & Yes (3) \\
Q18 & 0.0036 & 0.8208 & No (0) & Yes (3) \\
Q19 & 0.0032 & 0.6728 & No (0) & Yes (1) \\
Q20 & 0.0031 & 0.6511 & No (0) & Yes (1) \\
Q21 & 0.0039 & 0.6982 & No (0) & Yes (1) \\
Q22 & 0.0034 & 0.6865 & No (0) & Yes (1) \\
Q23 & 0.0035 & 0.6964 & No (0) & Yes (1) \\
Q24 & 0.0031 & 0.6937 & No (0) & Yes (1) \\
Q25 & 0.0033 & 0.6759 & No (0) & Yes (1) \\
\midrule
\multicolumn{5}{l}{\textbf{Summary Retrieval Geometry Metrics:}} \\
\multicolumn{5}{l}{$\bullet$ Embedding Norm: Pre $\|\mathbf{v}\| \approx 19.8$ ($d^2 \approx 280$) $\to$ Post $\|\mathbf{v}\| = 1.0$ (Unit-$L_2$)} \\
\multicolumn{5}{l}{$\bullet$ Total Retrieval Events: Pre = $0$ / 25 queries $\to$ Post = $16$ retrieval events} \\
\multicolumn{5}{l}{$\bullet$ Distinct Retrieval Queries: Pre = $0$ ($0.0\%$) $\to$ Post = $12$ of 25 ($48.0\%$)} \\
\bottomrule
\multicolumn{5}{@{}p{0.90\linewidth}@{}}{\scriptsize \textsuperscript{*}Retrieval threshold $\tau=0.50$ for all queries; a retrieval event is counted when $S\geq\tau$. Primary sources: \url{benchmark/pre_remediation_results.csv}, \url{benchmark/seed42/benchmark_results.csv}, and \url{manuscript/figures/source/fig3_retrieval_geometry.py}. Query-level similarity scores use $S=1/(1+d^2)$. The pre-remediation condition exhibits $S\approx0.0035\ll\tau$, whereas Unit-$L_2$ normalization increases the observed similarity for retrieved contexts.} \\
\end{tabular}
\end{table*}

\begin{table*}[t]
\centering
\caption{Numerical Data Underlying the PostgreSQL Execution Success vs. Relational Execution Accuracy Analysis Across the Six Experimental Modes}
\label{tab:fig4_validation}
\footnotesize
\setlength{\tabcolsep}{8pt}
\renewcommand{\arraystretch}{1.06}
\begin{tabular}{l c c c c c c}
\toprule
\textbf{Experimental Mode} & \textbf{ExecSucc (\%)} & \textbf{ExecAcc (\%)} & \textbf{Gap (pp)} & \textbf{PG Successes / 75} & \textbf{Rel. Correct / 75} & \textbf{Total Evals} \\
\midrule
Mode 1 (Zero-Shot Baseline) & $76.00 \pm 0.00$ & $57.33 \pm 2.31$ & $18.67$ & 57 / 75 & 43 / 75 & 75 \\
Mode 2 (Stateless Self-Correction) & $92.00 \pm 0.00$ & $68.00 \pm 0.00$ & $24.00$ & 69 / 75 & 51 / 75 & 75 \\
Mode 3 (Naive Vector RAG) & $92.00 \pm 0.00$ & $68.00 \pm 0.00$ & $24.00$ & 69 / 75 & 51 / 75 & 75 \\
Mode 4 (Full ARMG) & $96.00 \pm 0.00$ & $68.00 \pm 0.00$ & $28.00$ & 72 / 75 & 51 / 75 & 75 \\
Mode 5 (ARMG $-$ Neg Constraints) & $96.00 \pm 0.00$ & $68.00 \pm 0.00$ & $28.00$ & 72 / 75 & 51 / 75 & 75 \\
Mode 6 (ARMG with $\lambda = 0.0$) & $94.67 \pm 2.31$ & $68.00 \pm 0.00$ & $26.67$ & 71 / 75 & 51 / 75 & 75 \\
\midrule
\textbf{Total Corpus Execution} & $91.11$ & $66.22$ & $24.89$ & 410 / 450 & 298 / 450 & 450 \\
\bottomrule
\multicolumn{7}{@{}p{0.92\linewidth}@{}}{\footnotesize \textsuperscript{*}Values are mean $\pm$ sample standard deviation across three isolated repeated executions (seeds 42, 123, and 999). Each mode contains 75 query evaluations.} \\
\end{tabular}
\end{table*}

\begin{table*}[t]
\centering
\caption{Numerical Data Underlying the Operational Trade-Off Profile: Full ARMG Relative to the Stateless Self-Correction Baseline}
\label{tab:fig5_validation}
\footnotesize
\setlength{\tabcolsep}{10pt}
\renewcommand{\arraystretch}{1.06}
\begin{tabular}{l c c c c c}
\toprule
\textbf{Evaluation Metric} & \textbf{Mode 2} & \textbf{Mode 4} & \textbf{Observed Delta} & \textbf{Delta Type} & \textbf{Unit} \\
\midrule
Mean Repair Iterations & $0.43$ & $0.28$ & $-34.38\%$ & Relative & retries/query \\
Mean Token Expenditure & $572.72$ & $542.37$ & $-5.30\%$ & Relative & tokens/query \\
PostgreSQL Execution Success & $92.00\%$ & $96.00\%$ & $+4.00\text{ pp}$ & Percentage points & pp \\
End-to-End Latency & $7,149.68$ & $8,564.89$ & $+19.79\%$ & Relative & ms \\
Relational Execution Accuracy & $68.00\%$ & $68.00\%$ & $0.00\text{ pp}$ & Percentage points & pp \\
\bottomrule
\multicolumn{6}{@{}p{0.88\linewidth}@{}}{\scriptsize \textsuperscript{*}Rate-based operational metrics (retries, tokens, and latency) are reported as relative percentage changes. Execution success and relational execution accuracy are reported as percentage-point differences.} \\
\end{tabular}
\end{table*}

\begin{table*}[t]
\centering
\caption{Numerical Data Underlying Persistent Memory Store Growth: Store Size and Governance Events Across the 25-Query Benchmark}
\label{tab:fig6_validation}
\footnotesize
\setlength{\tabcolsep}{12pt}
\renewcommand{\arraystretch}{1.03}
\begin{tabular}{l c c l l}
\toprule
\textbf{Query} & \textbf{Mode 3 Store Size} & \textbf{Mode 4 Store Size} & \textbf{Mode 3 Event} & \textbf{Mode 4 Governance Event} \\
\midrule
Q01 & 1 & 0 & \texttt{NAIVE\_STORED} & None (Att.~1 Succ) \\
Q02 & 2 & 0 & \texttt{NAIVE\_STORED} & None (Att.~1 Succ) \\
Q03 & 3 & 0 & \texttt{NAIVE\_STORED} & None (Att.~1 Succ) \\
Q04 & 3 & 1 & None (Att.~1 Fail) & \texttt{ADMITTED} \\
Q05 & 4 & 1 & \texttt{NAIVE\_STORED} & None (Att.~1 Succ) \\
Q06 & 5 & 1 & \texttt{NAIVE\_STORED} & None (Att.~1 Succ) \\
Q07 & 6 & 1 & \texttt{NAIVE\_STORED} & None (Att.~1 Succ) \\
Q08 & 7 & 1 & \texttt{NAIVE\_STORED} & None (Att.~1 Succ) \\
Q09 & 8 & 1 & \texttt{NAIVE\_STORED} & None (Att.~1 Succ) \\
Q10 & 9 & 1 & \texttt{NAIVE\_STORED} & None (Att.~1 Succ) \\
Q11 & 10 & 1 & \texttt{NAIVE\_STORED} & \texttt{REJECTED} (Fail) \\
Q12 & 11 & 1 & \texttt{NAIVE\_STORED} & None (Att.~1 Succ) \\
Q13 & 12 & 2 & \texttt{NAIVE\_STORED} & \texttt{ADMITTED} \\
Q14 & 13 & 2 & \texttt{NAIVE\_STORED} & None (Att.~1 Succ) \\
Q15 & 13 & 3 & None (Att.~1 Fail) & \texttt{ADMITTED} \\
Q16 & 14 & 3 & \texttt{NAIVE\_STORED} & None (Att.~1 Succ) \\
Q17 & 15 & 3 & \texttt{NAIVE\_STORED} & \texttt{REINFORCED} (mem-Q13) \\
Q18 & 16 & 3 & \texttt{NAIVE\_STORED} & None (Att.~1 Succ) \\
Q19 & 17 & 3 & \texttt{NAIVE\_STORED} & None (Att.~1 Succ) \\
Q20 & 18 & 3 & \texttt{NAIVE\_STORED} & None (Att.~1 Succ) \\
Q21 & 19 & 3 & \texttt{NAIVE\_STORED} & None (Att.~1 Succ) \\
Q22 & 20 & 3 & \texttt{NAIVE\_STORED} & None (Att.~1 Succ) \\
Q23 & 21 & 3 & \texttt{NAIVE\_STORED} & None (Att.~1 Succ) \\
Q24 & 22 & 3 & \texttt{NAIVE\_STORED} & None (Att.~1 Succ) \\
Q25 & 23 & 3 & \texttt{NAIVE\_STORED} & None (Att.~1 Succ) \\
\midrule
\multicolumn{5}{l}{\textbf{Key Governance Invariants:}} \\
\multicolumn{5}{l}{$\bullet$ Initial Admissions: Q04 (Store=1), Q13 (Store=2), Q15 (Store=3)} \\
\multicolumn{5}{l}{$\bullet$ Mutual Exclusion Invariant: Q17 reinforces \texttt{mem-Q13}; new admission suppressed} \\
\multicolumn{5}{l}{$\bullet$ Observed Mode 4 Store Plateau: Store size remains at 3 from Q15 through Q25 in the evaluated sequential benchmark} \\
\multicolumn{5}{l}{$\bullet$ Mode 3 Unmanaged Accumulation: 23 entries accumulated without lifecycle governance} \\
\bottomrule
\multicolumn{5}{@{}p{0.92\linewidth}@{}}{\scriptsize \textsuperscript{*}Primary sources: \url{benchmark/seed42/benchmark_results.csv}, \url{benchmark/seed123/benchmark_results.csv}, \url{benchmark/seed999/benchmark_results.csv}, and \url{manuscript/figures/source/fig6_memory_growth.py}. Progression is identical across all three benchmark seeds ($n = 3$). In Mode 3, queries failing on Attempt 1 without self-correction (Q04, Q15) are unadmitted, while 23 successful queries are stored unconstrained. In Mode 4, only post-repair recoveries passing admission gating ($\text{Utility}_0 \ge 0.25$) are admitted, reaching a persistent plateau of exactly 3 memories.} \\
\end{tabular}
\end{table*}

\begin{table*}[t]
\centering
\caption{Traceability Matrix Linking Publication Diagrams to Underlying Numerical Validation Tables and Raw Data Sources}
\label{tab:figure_validation_matrix}
\footnotesize
\setlength{\tabcolsep}{4pt}
\renewcommand{\arraystretch}{1.08}
\begin{tabularx}{\textwidth}{@{}p{4.2cm} p{2.1cm} p{4.2cm} p{2.2cm} X@{}}
\toprule
\textbf{Diagram / Chart Title} & \textbf{Validation Table} & \textbf{Primary Source Dataset} & \textbf{Data Rows} & \textbf{Reproducible From Source?} \\
\midrule
FAISS Retrieval Geometry & Table~\ref{tab:fig3_validation} & \url{benchmark/pre_remediation_results.csv} & 25 queries (Q01--Q25) & Reproducible from source \\
PostgreSQL Execution Success vs. Relational Execution Accuracy & Table~\ref{tab:fig4_validation} & \url{benchmark/seed*/benchmark_results.csv} & 6 experimental modes & Reproducible from source \\
Operational Trade-Off Profile & Table~\ref{tab:fig5_validation} & \url{benchmark/seed*/benchmark_results.csv} & 5 operational metrics & Reproducible from source \\
Persistent Memory Store Growth & Table~\ref{tab:fig6_validation} & \url{benchmark/seed*/benchmark_results.csv} & 25 sequential queries & Reproducible from source \\
\bottomrule
\multicolumn{5}{@{}p{\linewidth}@{}}{\scriptsize All validation tables are derived from the repository benchmark data and figure-generation scripts and were checked using the automated validation suite.} \\
\end{tabularx}
\end{table*}

\bibliographystyle{IEEEtran}
\bibliography{\jobname}
\end{document}